\documentclass[12pt]{article}

% ============================================================
% PACKAGES
% ============================================================

\usepackage{amsmath}      % Mathematical environments: align, equation, etc.
\usepackage{amsthm}       % Theorem, definition, proof environments
\usepackage{amsfonts}     % Mathematical fonts such as \mathbb
\usepackage{amssymb}      % Additional mathematical symbols
\usepackage{mathtools}    % Extensions to amsmath
\usepackage{geometry}     % Page margins
\usepackage{enumitem}     % Customize enumerate/itemize
\usepackage{mdframed}     % Framed boxes
\usepackage{tikz}         % Draw diagrams
\usepackage{hyperref}     % Clickable references and links

\geometry{
    letterpaper,
    margin=0.8in
}

\linespread{1.2}


% ============================================================
% NUMBER SYSTEMS
% ============================================================
% Instead of writing \mathbb{R} every time, use \RR.

\newcommand{\RR}{\mathbb{R}}   % Real numbers
\newcommand{\CC}{\mathbb{C}}   % Complex numbers
\newcommand{\NN}{\mathbb{N}}   % Natural numbers
\newcommand{\ZZ}{\mathbb{Z}}   % Integers
\newcommand{\QQ}{\mathbb{Q}}   % Rational numbers


% ============================================================
% COMMON MATHEMATICAL NOTATION
% ============================================================
% These commands are useful for frequently used notation.

\newcommand{\set}[1]{\left\{#1\right\}}
% Example:
%   \set{x \in \RR : x > 0}
% produces:
%   { x in R : x > 0 }

\newcommand{\abs}[1]{\left|#1\right|}
% Example:
%   \abs{x-y}
% produces:
%   |x-y|

\newcommand{\norm}[1]{\left\|#1\right\|}
% Example:
%   \norm{x}
% produces:
%   ||x||

\newcommand{\inner}[2]{\left\langle #1,#2\right\rangle}
% Example:
%   \inner{x}{y}
% produces:
%   <x,y>

\newcommand{\parens}[1]{\left(#1\right)}
% Example:
%   \parens{x+y}

\newcommand{\brac}[1]{\left[#1\right]}
% Example:
%   \brac{a,b}

\newcommand{\angles}[1]{\left\langle #1\right\rangle}
% Example:
%   \angles{x,y}


% ============================================================
% THEOREM / THEORY COMPONENTS
% ============================================================

% "definition" style is upright text.
% Use this for definitions, examples, and remarks.

\theoremstyle{definition}

% ------------------------------------------------------------
% DEFINITION
% ------------------------------------------------------------
% Use when introducing a formal mathematical definition.
%
% Usage:
%
% \begin{definition}[Metric]
% A metric on a set X is a function ...
% \end{definition}
%
% [Metric] is optional and gives the definition a name.

\newtheorem{definition}{Definition}[section]


% ------------------------------------------------------------
% EXAMPLE
% ------------------------------------------------------------
% Use to give a concrete example of a definition/theorem.
%
% Usage:
%
% \begin{example}
% Consider X = \RR with the usual distance.
% \end{example}

\newtheorem{example}[definition]{Example}


% ------------------------------------------------------------
% REMARK
% ------------------------------------------------------------
% Use for observations, comments, or small additional facts.
% It is NOT intended for major mathematical results.
%
% Usage:
%
% \begin{remark}
% This definition does not require X to be a subset of \RR^n.
% \end{remark}

\newtheorem{remark}[definition]{Remark}


% ============================================================
% THEOREM-LIKE COMPONENTS
% ============================================================

% "plain" style makes the mathematical statements italic.

\theoremstyle{plain}


% ------------------------------------------------------------
% THEOREM
% ------------------------------------------------------------
% Use for an important/general mathematical result.
%
% Usage:
%
% \begin{theorem}[Cauchy--Schwarz Inequality]
% For all x,y \in \RR^n,
% \[
% |\inner{x}{y}| \leq \norm{x}\norm{y}.
% \]
% \end{theorem}

\newtheorem{theorem}[definition]{Theorem}


% ------------------------------------------------------------
% PROPOSITION
% ------------------------------------------------------------
% Use for a useful mathematical result that is usually less
% central than a theorem.
%
% Usage:
%
% \begin{proposition}
% Every finite-dimensional subspace of \RR^n is closed.
% \end{proposition}

\newtheorem{proposition}[definition]{Proposition}


% ------------------------------------------------------------
% LEMMA
% ------------------------------------------------------------
% Use for an auxiliary result whose main purpose is to help
% prove another theorem/proposition.
%
% Usage:
%
% \begin{lemma}
% ...
% \end{lemma}

\newtheorem{lemma}[definition]{Lemma}




% ------------------------------------------------------------
% COROLLARY
% ------------------------------------------------------------
% Use for a result that follows relatively directly from a
% theorem/proposition that was just established.
%
% Usage:
%
% \begin{corollary}
% ...
% \end{corollary}

\newtheorem{corollary}[definition]{Corollary}


% ============================================================
% PROOF
% ============================================================
% The proof environment comes from amsthm.
%
% Usage:
%
% \begin{proof}
% Suppose that ...
%
% Therefore, ...
% \end{proof}
%
% LaTeX automatically adds the QED symbol at the end.


% ============================================================
% CUSTOM BOX: IDEA
% ============================================================
% Use this when you want to record intuition or the main idea
% behind a mathematical concept.
%
% Usage:
%
% \begin{idea}
% \textbf{Key idea.}
% A metric gives us a notion of distance...
% \end{idea}

\newmdenv[
    linewidth=1pt,
    skipabove=10pt,
    skipbelow=10pt
]{idea}


% ============================================================
% CUSTOM BOX: WARNING
% ============================================================
% Use this for common mistakes, distinctions, or things that
% are easy to confuse.
%
% Usage:
%
% \begin{warning}
% Do not confuse continuity with uniform continuity.
% \end{warning}

\newmdenv[
    linewidth=1pt,
    skipabove=10pt,
    skipbelow=10pt
]{warning}

% ============================================================
% CUSTOM BOX: NOTE
% ============================================================
% Use this for a short personal note, useful observation,
% connection between concepts, or explanation.
%
% Unlike "Remark", this is intended for informal notes.
%
% Usage:
%
% \begin{note}
% A sequence can be viewed as a function from $\NN$ to $\RR$.
% \end{note}

\newmdenv[
    linewidth=1pt,
    skipabove=10pt,
    skipbelow=10pt
]{note}


% ============================================================
% DOCUMENT
% ============================================================

\begin{document}


% ============================================================
% TITLE
% ============================================================

\begin{center}
    {\Large \textbf{Cauchy sequences, Complete spaces and limsup}}\\[0.5em]
    {\large Lecture 17}
\end{center}

\tableofcontents

\newpage

\section{Cauchy sequences}

\begin{definition}[Monotonically increasing sequence]
$\{s_n\}$ is monotonically increasing if $s_n \le s_{n+1}$ for all $n$.
\end{definition}

\begin{theorem}[Convergence of bounded and monotonic sequences]
Bounded and monotonic (increasing or decreasing) sequences converge (to its $\inf$ or $\sup$).
\end{theorem}

\begin{definition}[Cauchy sequences]
$\{p_n\}$ is Cauchy if for all $\epsilon > 0$, there is a $N$ such that $m,n > N$ then $d(p_m,p_n)<\epsilon$.
\end{definition}

\begin{note}
past index $N$, every points are $\epsilon$-close.
\end{note}

\begin{theorem}[Cauchy-ness of converging sequences]
If $\{s_n\}$ converges, then $\{s_n\}$ is Cauchy.
\end{theorem}

\begin{note}
No inverse direction. Counter example, if you take a sequence in $\QQ$ that is Cauchy (converge to $\pi$, for example), but the sequence doesn't converge in $\QQ$.
\end{note}

\section{Complete spaces}

\begin{definition}[Complete space]
A metric space $X$ is complete if every Cauchy sequences in $X$ converge (to a point of $X$).
\end{definition}

\begin{theorem}[Some complete spaces]
Compact spaces, closed subsets of complete space, $\RR^k$ are complete.
\end{theorem}

\begin{theorem}[Very strong theorem]
Every metric space $(X, d)$ has a completion $(X^*, \Delta)$. We call $X^*$ extends $X$ and $\Delta$ extends $d$ or $X$ is isometrically embedded in $X^*$ and $X^*$ is complete (extra credit).
\end{theorem}

\section{limsup}
\begin{definition}[limsup and liminf]
For sequence $\{s_n\}$ in $\RR$, let $E = \{x: x \text{ is a subsequence's limit}\}$ (may include $\infty$). Let $s^*= \sup E = \limsup \{s_n\}$ and $s_*= \inf E = \liminf \{s_n\}$.
\end{definition}

\begin{note}
There is another definition beside the definition above, which is $\limsup s_n = \lim_{n\rightarrow\infty}(\sup_{k\ge n}s_k)$. This is equivalent to the definition above, but it is very convenient for computing and proving (maybe). But the more intuitive is the definition above, which is the upper limit.
\end{note}

\begin{example}
$s_n =\{.1, \frac{3}{2}, .11, \frac{4}{3},.111,\frac{5}{4},...\}$ with $\limsup s_n = 1$, $\liminf s_n = \frac{1}{9}$.
\end{example}

\begin{theorem}[Existence of $\limsup$]
$s^* \in E$ (We can say that $\limsup$ always exists).
\end{theorem}

\begin{theorem}[Weird theorem of $\limsup$]
Given any $x > \limsup \{s_n\}$, then there is a $N$ that if $n \ge N$, then $s_n < x$.  
\end{theorem}

\begin{note}
In this theorem, try not to fall for the pitfall of thinking $\limsup$ is greater than any element in the original sequence. In original sequence, there's a point that is greater than $\limsup$ is very normal, $\limsup$ is just the greatest among those subsequences' limit.
\end{note}
\end{document}